\documentclass{mipt-thesis-ms}
% Следующие две строки нужны только для biblatex. Для inline-библиографии их следует убрать.
%\usepackage{mipt-thesis-biblatex}
%\addbibresource{example.bib}

\title{Разработка системы для разбора библиографических ссылок на основе дообученных больших языковых моделей с использованием LoRA-адаптеров}
\author{Красоткин С.\,А.}
\supervisor{Грабовой А.\,В.}
%\referee{Петров Д.\,Е.}       % требуется только для mipt-thesis-ms
\groupnum{М05-411е}
\faculty{Физтех-школа Прикладной Математики и Информатики (ФПМИ)}
\department{Кафедра корпоративных информационных систем}

\begin{document}

\frontmatter

\titlecontents

\chapter{Реферат}
Отчёт
% объект исследования или разработки;
% методы или методологию проведения работы;
% область применения результатов;
% рекомендации по внедрению или итоги внедрения результатов НИР;
% экономическую эффективность или значимость работы;
% прогнозные предположения о развитии объекта исследования.

% цель работы;
Цель работы: разработка системы для разбора библиографии на основе дообученных больших языковых моделей (БЯМ) с использованием LoRA-адаптеров.

% результаты работы и их новизну;
Результат работы: Сравнение проведено для двух уровней задачи: (i) детектирование позиций/границ библиографии и извлечение сырых сносок; (ii) структурирование ссылок по полям, включая классовое выделение авторов. По RedUp API (гистограммные градиентные бустинговые деревья и регулярные выражения для идентификации границ секций; \cite{kopanichuk}) точность определения позиций секции составила порядка 65\%, при этом для извлечения сырых сносок получены значения $F_1\approx0{.}86$ и точности $\approx0{.}80$; для классового выделения авторов $F_1\approx0{.}93$ при точности $\approx0{.}89$. Для дообученного GROBID (\cite{kiryanov}) на извлечении сырых ссылок зафиксированы $F_1\approx0{.}97$ и точность $\approx0{.}98$, тогда как для классового выделения авторов — $F_1\approx0{.}53$ и точность $\approx0{.}54$. Значения интерпретируются совместно с качественными примерами и зависят от стиля оформления и уровня шума источника.

\mainmatter


\chapter{Введение}

Актуальной задачей современной научной коммуникации является аккуратное извлечение и приведение к единому виду библиографических ссылок в статьях. Реальная практика многообразна: стили оформления различаются, источники содержат шум, а одна и та же сущность может описываться по‑разному. В работе мы рассматриваем связанный процесс из трёх шагов: сначала находится сама библиографическая секция в тексте, затем из неё выделяются отдельные ссылки, после чего каждая ссылка раскладывается по полям (авторы, заголовок, источник и др.). Цель исследования — сравнить применимость доступных подходов на разнородных корпусах: промышленного пайплайна RedUpAPI, специализированного инструмента GROBID и обобщённого LLM‑подхода. Показано, что RedUpAPI надёжен для обнаружения границ секции (точность около 65 процентов) и извлечения сырых ссылок (F1‑мера около 0.86 при точности около 0.80), а также даёт высокие результаты при классовом выделении авторов (F1‑мера около 0.93 при точности около 0.89). Дообученный GROBID демонстрирует очень высокое качество на извлечении сырых ссылок (F1‑мера около 0.97 при точности около 0.98), но заметно уступает при выделении авторов (F1‑мера около 0.53 при точности около 0.54). Практическая ценность результатов — поддержка «обратного библиографического списка» и интеграция в издательские и репозитории для регулярного автоматического обновления ссылок.

\chapter{Литературный обзор}

...

\chapter{Процесс исследований}

\section{Методы исследований}

\section{Экспериментальные работы}

\chapter{Обсуждение результатов}

\section{Оценку полноты}

\section{Оценку полноты}

\section{Дальнейшие работы}

\section{Оценка достоверности полученных результатов}

\chapter{Заключение}


\backmatter

\printbib
% Следующие строки необходимо раскомментировать, а предыдущую закомментировать, если используется inline-библиография.
%\begin{thebibliography}{99}
%    \bibitem{langmuir26}
%        H. Mott-Smith, I. Langmuir. ``The theory of collectors in gaseous discharges''. \emph{Phys. Rev.} \textbf{28} (1926)
%\end{thebibliography}

\end{document}
